% section: 式の変形

\section{式の変形}\label{節：式の変形}
  \noindent\textbf{【本編で使う場面】}\par
  同じ値を保ちながら既知の形を作る操作を確認する.定数の調整は\sectionref*{節：特性方程式型},\bookterm{partial}{部分分数分解}は\sectionref*{節：階差型}で和を計算する際に使う.

  \subsection{\texorpdfstring{$+a-a$をする}{+a-aをする}}\label{節：+a-aをする}
    二次式の中に、二乗の形を作る。
    $a\ne0$として、足りない定数を足し、同じだけ引くと、
    \begin{align*}
      &ax^2+bx+c\\
      &\qquad=a\left(x^2+\frac{b}{a}x\right) +c\\
      &\qquad=a\left(x^2+\frac{b}{a}x+\frac{b^2}{4a^2}-\frac{b^2}{4a^2}\right)+c \tag*{\(\cdots\) (☆)}\\
      &\qquad=a\left(x+\frac{b}{2a}\right)^2-\frac{b^2}{4a}+c
    \end{align*}
    となる。☆の行で、二乗に必要な定数を補った。
    足して引いた量は零なので、式の値は変わらない。
    $x$を含む部分が一つの二乗にまとまり、式の形が見やすくなる。
    この操作を\bookterm{square}{平方完成}という。
    二次関数の最大・最小や、軸の位置を調べるときにも使える。

    値を保ったまま、扱いやすい形を作る。
    同じ発想が、後の定数項を取り除く操作につながる。

  \subsection{部分分数分解}\label{節：部分分数分解}
    次は、積を含む分母を、より簡単な分数の和へ分ける。
    \[
      \frac{1}{(x+a)(x+b)}
    \]
    は、そのままでは和を計算しにくい。
    分母を一次の式に分けるため、次の形を作る。
    ここでは$x+a\ne0$、$x+b\ne0$とする。
    \[
      \frac{1}{(x+a)(x+b)}=\frac{\alpha}{x+a}+\frac{\beta}{x+b}
    \]
    $\alpha,\beta$について求めると,
    \begin{align*}
    \frac{1}{(x+a)(x+b)}
      &= \frac{\alpha}{x+a}+\frac{\beta}{x+b}\\
      &= \frac{(\alpha+\beta)x+\alpha b+\beta a}
      {(x+a)(x+b)}
    \end{align*}
    \[
    \therefore\quad
    \begin{cases}
    \alpha+\beta=0,\\
    \alpha b+\beta a=1
    \end{cases}
    \]
    これを解くと,\,$\beta=-\alpha$であるから,
    \[
      \alpha b-\alpha a=1
      \iff \alpha(b-a)=1
    \]
    したがって,\,$a\ne b$であれば,
    \[
      \alpha=\frac{1}{b-a},
      \qquad
      \beta=-\frac{1}{b-a}
    \]
    が得られる(\,$a=b$のときはそもそも分母が重解$(x+a)^2$になり,この形の分解はできない).
    二つの係数を定めると、一次の分母をもつ分数の和に分かれる。
    この操作を\bookterm{partial}{部分分数分解}という。
    3次の場合も,\,$p,q,r\ne0$で,分母の一次因子$px+a,qx+b,rx+c$の根
    \[
      -\frac{a}{p},\qquad -\frac{b}{q},\qquad -\frac{c}{r}
    \]
    が互いに異なるときは,
\BookMathLayout{\[
      \frac{1}{(px+a)(qx+b)(rx+c)}=\frac{\alpha}{px+a}+\frac{\beta}{qx+b}+\frac{\gamma}{rx+c}
    \]}{
\[
\begin{aligned}
\frac{1}{(px+a)(qx+b)(rx+c)}
&=\frac{\alpha}{px+a}+\frac{\beta}{qx+b}\\
&\quad+\frac{\gamma}{rx+c}
\end{aligned}
\]
}
    と置き,係数$\alpha,\beta,\gamma$を定めることで分解できる.
    言い換えれば,これらの一次因子が互いに定数倍でないことが条件である.
    重複因子があるときは,その重複度に応じて例えば$\dfrac{A}{x+a}+\dfrac{B}{(x+a)^2}$のような項も必要になる.
